% !TEX root = math-notes.tex
\documentclass[11pt,a4paper]{article}
\usepackage[margin=2.4cm]{geometry}
\usepackage{amsmath,amssymb,amsthm}
\usepackage{stmaryrd}
\usepackage{enumitem}
\usepackage{xcolor}
\usepackage[colorlinks,linkcolor=blue!50!black]{hyperref}

\theoremstyle{definition}
\newtheorem{definition}{Definition}
\theoremstyle{plain}
\newtheorem{lemma}{Lemma}
\newtheorem{theorem}{Theorem}
\newtheorem{corollary}{Corollary}
\theoremstyle{remark}
\newtheorem*{remark}{Remark}

\newcommand{\Dd}{\mathcal{D}}
\newcommand{\Tt}{\mathcal{T}}
\newcommand{\dataf}{\mathsf{data}}
\newcommand{\disc}{\mathsf{disc}}
\newcommand{\reach}{\mathrm{reach}}
\newcommand{\ev}[2]{\llbracket #1 \rrbracket_{#2}}

\title{Math toolkit for the delegation / coverage argument}
\author{Working notes}
\date{}

\begin{document}
\maketitle
\vspace{-2em}

\section{The two empty things}

Think of a SPARQL result as a \emph{table}.

\begin{center}
\begin{tabular}{lll}
\hline
Object & Notation & Meaning \\
\hline
Empty solution \emph{set}     & $\emptyset$              & no rows: the pattern did not match \\
Empty solution \emph{mapping} & $\mu_\emptyset$          & one row, zero columns \\
Join identity                 & $\{\mu_\emptyset\}$      & $\Omega \bowtie \{\mu_\emptyset\} = \Omega$ \\
\hline
\end{tabular}
\end{center}

\noindent
The cleanest instance is \textsf{ASK}, which is exactly what $\Delta$ computes:
\[
  \textsf{ASK}\ \{\, \langle a\rangle\ \langle p\rangle\ \langle b\rangle \,\}
  \;=\;
  \begin{cases}
    \textsf{true}  & \text{iff } \ev{\cdot}{G} = \{\mu_\emptyset\},\\
    \textsf{false} & \text{iff } \ev{\cdot}{G} = \emptyset.
  \end{cases}
\]
So $\{\mu_\emptyset\}$ means \emph{yes} and $\emptyset$ means \emph{no}. Note
$\{\mu_\emptyset\} \neq \emptyset$: it is a one-element set. Writing $\emptyset$ for
``matched with no bindings'' would invert the indicator. Notation: $\mu_\emptyset$ in
P\'erez--Arenas--Guti\'errez, $\mu_0$ in SPARQL~1.1 \S18.5, which also names
$Z = \{\mu_0\}$ the join identity.

\section{The setting as a parameter}
\label{sec:setting}

Nothing in what follows fixes what $\Dd$ \emph{is}. Every definition and both theorems
use only three properties: sources hold graphs, query evaluation is monotone in the
graph, and the reference answer is evaluation over the union of the sources. What varies
between federation and traversal is not the algebra but \emph{how the set of sources
comes to be known}, so we make that a parameter.

\begin{definition}[Discovery setting]
\label{def:setting}
A \emph{discovery setting} is a tuple
$W = \langle \Dd, \dataf, \Dd_0, \disc \rangle$: a set of sources $\Dd$, a map
$\dataf : \Dd \to 2^{\Tt}$ giving the graph each one holds, a set of seeds
$\Dd_0 \subseteq \Dd$, and a \emph{discovery function} $\disc$ assigning to every BGP $B$
a map $\disc_B : \Dd \times 2^{\Tt} \to 2^{\Dd}$, where $\disc_B(d, \dataf(d))$ are the
sources that a retrieved $d$ exposes while $B$ is evaluated.
\end{definition}

\begin{definition}[Reachable set, reference answer]
\label{def:reach}
Iterating discovery from the seeds, let
$\Dd_{i+1} := \Dd_i \cup \bigcup_{d \in \Dd_i} \disc_B\bigl(d, \dataf(d)\bigr)$. The
\emph{reachable set} is $\reach_B(W) := \bigcup_{i \geq 0} \Dd_i$, the \emph{reference
answer} to $B$ is $\ev{B}{G_{\reach_B(W)}}$, and correctness of a delegation is always
relative to it. Equivalently $\reach_B(W)$ is the smallest set containing $\Dd_0$ and
closed under $\disc_B$, though nothing below uses that characterisation.
\end{definition}

\noindent
The query is a genuine parameter of discovery --- $c_{\mathrm{All}}$ and
$c_{\mathrm{Match}}$ differ in nothing else --- so $\disc$ and $\reach$ carry it. A
single $B$ is fixed from here on and we drop the subscript, writing $\disc$ and
$\reach(W)$.

\begin{center}
\small\begin{tabular}{@{}p{5.1cm}p{2.7cm}l@{}}
\hline
$\disc$ & $\Dd_0$ & recovers \\
\hline
$\disc \equiv \emptyset$
  & all members
  & classical federated query processing \\
all URIs occurring in $\dataf(d)$
  & seed URIs
  & LTQP under $c_{\mathrm{All}}$ \\
URIs of triples matching some $tp \in B$
  & seed URIs
  & LTQP under $c_{\mathrm{Match}}$ \\
the subweb spec.\ inside $\dataf(d)$
  & the user's document
  & guided traversal \\
a link path expression from the query
  & seed URIs
  & LDQL \\
\hline
\end{tabular}
\end{center}

\begin{definition}[State at time $T$, frontier, termination]
\label{def:state}
A \emph{state} at time $T$ is a set $\Dd'(T)$ with
$\Dd_0 \subseteq \Dd'(T) \subseteq \reach(W)$, non-decreasing in $T$. Its
\emph{frontier} is
\[
  F(T) \;:=\; \Bigl( \bigcup_{d \in \Dd'(T)} \disc\bigl(d, \dataf(d)\bigr) \Bigr)
              \setminus \Dd'(T),
\]
the sources exposed by what has been retrieved but not yet retrieved themselves.
\emph{Discovery has terminated} at $T$ if $F(T) = \emptyset$, equivalently
$\Dd'(T) = \reach(W)$.
\end{definition}

\begin{remark}[Federation is the degenerate case]
Classical federation is $\disc \equiv \emptyset$ with $\Dd_0$ the configured member
list. Then $\reach(W) = \Dd_0$ and $F(0) = \emptyset$: discovery terminates
before the first request. This is exactly why exclusive groups are sound there, and why
the impossibility result below has nothing to say about it.
\end{remark}

\begin{remark}[The reference answer is not the Web]
$\ev{B}{G_{\reach(W)}}$ is evaluation over the union of the \emph{reachable}
sources, not over the Web. Taking $\disc$ to expose every URI and $\Dd$ to be every
document would give full-web semantics, which is not finitely computable; the traversal
instantiations keep $\reach(W)$ bounded by the seeds and the criterion.
\end{remark}

\begin{remark}[How $\Delta$ is obtained differs]
In a federated setting $\Delta$ comes from cheap probes --- \textsf{ASK} requests, a VoID
description, an index --- independently of retrieving the data. Under traversal there is
no such probe: $\Delta(t,d)$ becomes known exactly when $\dataf(d)$ does, so $\Delta$
collapses into the state. The source-description language is therefore strictly poorer in
the traversal instantiations, which strengthens rather than weakens the impossibility
argument.
\end{remark}

\begin{remark}[Beyond RDF]
Removing triple patterns and dereferencing leaves monotone conjunctive queries over
horizontally partitioned data, where a fragment may be shipped whole to one partition
exactly when no answer spans partitions. That is the classical localisation condition for
distributed query processing; the RDF instance merely supplies particular notions of
partition and of link.
\end{remark}

\section{Relevance and delegation}

\noindent
\textbf{Notation.} For a set of sources $A \subseteq \Dd$ we write
$G_A := \bigcup_{d \in A} \dataf(d)$ for the union of their graphs. The three instances
used below are $G_{\Dd}$ (the whole federation), $G_{\Dd \setminus S}$ (the
non-delegated sources) and $G_A$ for an arbitrary cover $A$ in the coverage section.


\begin{definition}[Source relevance]
Let $\mathcal{P}$ be the set of triple patterns and $\Dd$ the set of all sources. Write
$\ev{t}{d}$ for the evaluation of $t$ over the data of $d$. The \emph{relevance
indicator} $\Delta : \mathcal{P} \times \Dd \to \{0,1\}$ is
\[
  \Delta(t,d) =
  \begin{cases}
    1, & \text{if } \ev{t}{d} \neq \emptyset,\\[2pt]
    0, & \text{otherwise,}
  \end{cases}
  \qquad
  \Dd_t := \{\, d \in \Dd \mid \Delta(t,d) = 1 \,\}.
\]
$\Delta$ is obtained for the patterns of the query under evaluation only --- one
\textsf{ASK} probe per pattern per source, as in FedX.
\end{definition}

\begin{definition}[Federation instance]
A \emph{federation instance} is a function $\delta : \Dd \to 2^{\Tt}$ giving the graph
held by each source. It is \emph{consistent} with $\Delta$ for $t_1,\dots,t_n$, written
$\delta \models \Delta$, if
\[
  \Delta(t_i,d) = 1 \iff \ev{t_i}{\delta(d)} \neq \emptyset
  \qquad \text{for all } i \in \{1,\dots,n\},\ d \in \Dd .
\]
\end{definition}

\begin{remark}
$\Delta$ is what the engine \emph{knows}; $\delta$ is what is \emph{there}. Many $\delta$
give the same $\Delta$, and the planner must be correct for all of them. Do not call this
a ``source assignment'': FedQPL Def.~8 already uses that term for a query plan.
\end{remark}

\begin{definition}[Delegation]
Let $B$ be a BGP and $B' \subseteq B$ the delegated fragment. Let $S \subseteq \Dd$ be a
set of sources such that every $d \in S$ provides an interface $I_d = (L_d, \varrho_d)$
with $B' \in L_d$; that is, each delegated source accepts $B'$ as a \emph{single}
request, returning $\varrho_d\big(B',\dataf(d)\big) = \ev{B'}{\dataf(d)}$. The delegated
plan obtains for $B'$
\[
  R \;:=\; \Big( \bigcup_{d \in S} \varrho_d\big(B', \dataf(d)\big) \Big)
           \;\cup\; \ev{B'}{G_{\Dd \setminus S}},
\]
and the delegation is \emph{correct} if
\[
  R \;\bowtie\; \ev{B \setminus B'}{G_{\Dd}} \;=\; \ev{B}{G_{\Dd}} .
\]
\end{definition}

\begin{remark}[Where the request granularity lives]
Semantically $\ev{B'}{\dataf(d)} = \bowtie_{t \in B'} \ev{t}{\dataf(d)}$, so over a
single source one conjunctive request and $|B'|$ pattern requests return the same
answers. The two \emph{plans} are distinguished by the order of union and join:
\[
  \bigcup_{d \in S} \ev{B'}{\dataf(d)}
  \qquad\text{versus}\qquad
  \bowtie_{t \in B'} \ \bigcup_{d \in S} \ev{t}{\dataf(d)} .
\]
Union outside the join is one request per source; join outside the union is one request
per pattern per source. Only the first is delegation.
\end{remark}

\begin{remark}[Capability and cost]
The condition $B' \in L_d$ is what makes ``single request'' meaningful, and it is the
point at which interface heterogeneity enters: a triple-pattern-only interface has
$L_d = \mathcal{P}$ and admits no $B'$ with $|B'| > 1$, while a composite resource
advertising the path $tp_1/tp_2$ has $L_d = \{tp_1/tp_2\}$, so $B' \in L_d$ is exactly
the requirement that the delegated fragment match what the resource publishes. Counting
requests as in FedQPL's sa-cost, delegation costs $|S|$ where the per-pattern route costs
$|B'| \cdot |S|$.
\end{remark}

\begin{remark}[Soundness is free]
$\dataf(d) \subseteq G_{\Dd}$ and $G_{\Dd \setminus S} \subseteq G_{\Dd}$, so
$R \subseteq \ev{B'}{G_{\Dd}}$; with $\bowtie$ monotone, the left-hand side is always a
subset of the right. Delegation can never produce a wrong answer, only lose results. The
whole content of the definition is completeness.
\end{remark}

\begin{remark}[Sanity checks]
$S = \emptyset$ gives $R = \ev{B'}{G_{\Dd}}$: not delegating is trivially correct.
$S = \Dd$ gives $R = \bigcup_{d \in \Dd} \ev{B'}{\dataf(d)}$: the no-inter-source-join
condition. $S = \{k\}$ with $k$ exclusive gives $R = \ev{B'}{\dataf(k)} =
\ev{B'}{G_{\Dd}}$: FedX's exclusive group, and FedQPL Def.~7 with
$\varphi = \mathrm{req}^{B'}_{k}$.
\end{remark}

\begin{remark}[What is lost]
Since $\ev{B}{G} = \ev{B'}{G} \bowtie \ev{B \setminus B'}{G}$, the gap between $R$ and
$\ev{B'}{G_{\Dd}}$ consists exactly of the $B'$-solutions spanning a boundary introduced
by the grouping: across the $S$ / $\Dd \setminus S$ split, or between two distinct
sources inside $S$. Such a solution is harmless when it has no join partner in
$B \setminus B'$, since it would have been discarded anyway.
\end{remark}

\begin{remark}[Why the two-block form does not work]
The variant $\ev{B}{\Dd} = \ev{B}{\Dd \setminus \{d\}} \cup \ev{B}{d}$ forbids only
joins crossing the $d$ boundary; it is blind to joins \emph{inside} the block
$\Dd \setminus \{d\}$. Take $\Dd = \{d, d_1, d_2\}$ and $B = \{t_1,t_2\}$ sharing a
variable, with $t_1$ matching only at $d_1$, $t_2$ only at $d_2$, and the two joining.
The equation holds, yet an inter-source join exists.
\end{remark}

\begin{theorem}[Correct delegation]
Let $B' = \{t_1,\dots,t_n\}$. Delegation of $B'$ to $S$ is correct for every federation
instance consistent with $\Delta$ if and only if
\[
  \Big( \exists\, k \in S : \bigcup_{i=1}^{n} \Dd_{t_i} \subseteq \{k\} \Big)
  \quad\text{or}\quad
  \Big( \bigcup_{i=1}^{n} \Dd_{t_i} \Big) \cap S = \emptyset .
\]
\end{theorem}

\begin{remark}
All relevant sources lie inside a single delegated source, or none of them lies in $S$
at all --- the latter being the ``delegating to an irrelevant source is harmless'' case,
which plain exclusivity does not capture. Anything in between admits a spanning witness.
Both disjuncts are forced by the poverty of $\Delta$: recording only per-pattern
relevance, it cannot reveal whether a join actually crosses sources. The quantifier over
instances is essential --- on a single fixed instance the ``only if'' fails, since two
sources holding identical data are both relevant yet delegating to either is correct.
\end{remark}

\section{Incomplete discovery}
\label{sec:incomplete}

Definition~1 gives $\Delta$ over all of $\Dd$, so an engine holding $\Delta$ knows
everything. To say anything about discovery we must separate what is true from what has
been seen.

\begin{definition}[Observation, completion, decision procedure]
\label{def:obs}
At time $T$ the engine has retrieved $\Dd'(T)$ and knows $\dataf(e)$ for each
$e \in \Dd'(T)$; write $O(T) := \bigl( \Dd'(T), (\dataf(e))_{e \in \Dd'(T)} \bigr)$. It
can therefore also compute the frontier $F(T)$, but knows nothing of $\dataf(d)$ for
$d \in F(T)$. A \emph{completion} of $O(T)$ is any discovery setting
$\widehat{W} = \langle \widehat{\Dd}, \widehat{\dataf}, \Dd_0, \disc \rangle$ agreeing
with $O(T)$ on $\Dd'(T)$; the adversary's freedom is exactly the data behind the
frontier. A \emph{decision procedure} is a partial function $\Pi$ from observations to
pairs $(B', S)$ with $S \subseteq \Dd'(T)$, and is \emph{sound} if whenever
$\Pi(O(T)) = (B', S)$, delegating $B'$ to $S$ is correct in \emph{every} completion
of $O(T)$.
\end{definition}

\begin{theorem}[No certificate under incomplete discovery]
\label{thm:nocert}
Let $O(T)$ be an observation at which discovery has not terminated, i.e.\ $F(T) \neq \emptyset$. Then there
are $B'$, $S \subseteq \Dd'(T)$ and completions $\widehat{\Dd}_1, \widehat{\Dd}_2$ of
$O(T)$ such that delegating $B'$ to $S$ is correct in $\widehat{\Dd}_1$ and incorrect in
$\widehat{\Dd}_2$. Consequently no sound decision procedure can output $(B', S)$.
\end{theorem}

\begin{proof}
Take $B' = B = \{t_1, t_2\}$ with
$t_1 = (?x, p, ?y)$ and $t_2 = (?y, q, ?z)$, so that $B \setminus B' = \emptyset$ and the
outer join in Definition~3 is with $\{\mu_\emptyset\}$, i.e.\ the identity. Let the
observation be $\Dd'(T) = \{e\}$ with $\dataf(e) = \{(a,p,b)\}$, and let $S = \{e\}$.

\emph{World 1.} Let $\widehat{\Dd}_1 = \{e\}$. Here $\Dd_{t_1} = \{e\}$ and
$\Dd_{t_2} = \emptyset$, so $\ev{B'}{G_{\widehat{\Dd}_1}} = \emptyset$, and likewise
$R = \emptyset$. The delegation is correct.

\emph{World 2.} Let $\widehat{\Dd}_2 = \{e, d^{*}\}$ with
$\dataf(d^{*}) = \{(b,q,c)\}$, where $d^{*} \in F(T)$ sits on the frontier. Now
\[
  \ev{B'}{G_{\widehat{\Dd}_2}} = \bigl\{\, \{?x \mapsto a,\ ?y \mapsto b,\ ?z \mapsto c\} \,\bigr\}
  \;\neq\; \emptyset,
\]
whereas both terms of $R$ vanish:
\[
  \bigcup_{d \in S} \ev{B'}{\dataf(d)} = \ev{B'}{\dataf(e)} = \emptyset,
  \qquad
  \ev{B'}{G_{\widehat{\Dd}_2 \setminus S}} = \ev{B'}{\dataf(d^{*})} = \emptyset,
\]
since $e$ has no $t_2$-match and $d^{*}$ has no $t_1$-match. Hence $R = \emptyset$ and
the answer is lost: the delegation is incorrect.

Both completions agree with $O(T)$ on $\Dd'(T)$, so no function of $O(T)$ can
distinguish them.
\end{proof}

\begin{corollary}[Infinite web]
\label{cor:infinite}
If $\reach(W)$ is infinite then $F(T) \neq \emptyset$ at every time $T$, and hence no
sound decision procedure can certify a delegation at any time.
\end{corollary}

\begin{proof}
At any time $T$ the engine has performed finitely many retrievals, so $\Dd'(T)$ is
finite. Suppose $F(T) = \emptyset$. Then $\disc\bigl(d, \dataf(d)\bigr) \subseteq \Dd'(T)$
for every $d \in \Dd'(T)$, i.e.\ $\Dd'(T)$ is closed under discovery, and
$\Dd_0 \subseteq \Dd'(T)$ by Definition~\ref{def:state}. By induction on $i$ we get
$\Dd_i \subseteq \Dd'(T)$ for every $i$: the base case is the seed inclusion, and if
$\Dd_i \subseteq \Dd'(T)$ then
$\Dd_{i+1} = \Dd_i \cup \bigcup_{d \in \Dd_i} \disc\bigl(d, \dataf(d)\bigr) \subseteq
\Dd'(T)$ by closure. Hence $\reach(W) = \bigcup_i \Dd_i \subseteq \Dd'(T)$ is finite,
contradicting the assumption. So $F(T) \neq \emptyset$, and Theorem~\ref{thm:nocert}
applies at every $T$.
\end{proof}

\begin{remark}[Never, not merely hard]
The corollary is stronger than a statement about cost. It is not that certification takes
a long time on a large web; it is that at no finite time does a certificate exist. Infinite
reachable sets are not exotic --- a server generating a document per natural number, or
per geographic area, suffices --- so this is the regime the traversal instantiations
actually live in. Contrapositively, certifying a delegation requires $\reach(W)$ to be
finite \emph{and} to have been retrieved in full.
\end{remark}

\begin{remark}[Why $\Delta$ alone cannot establish this]
The witness must exhibit concrete triples. At the level of $\Delta$ the two worlds are
distinguished only by $\Delta(t_2, d^{*}) = 1$, and that alone does not settle
correctness: if $d^{*}$ merely \emph{duplicated} data already held by $e$, no new
solution would appear and the delegation would remain correct. What breaks it is a
triple that is new \emph{and} joins across the boundary.
\end{remark}

\begin{remark}[Relevance elsewhere is not by itself a problem]
Discovering $\Delta(t, d) = 1$ for an undiscovered $d \notin S$ does \emph{not} break the
delegation. Since $d \in \Dd \setminus S$, the second term $\ev{B'}{G_{\Dd \setminus S}}$
of $R$ still collects everything $d$ can answer on its own. Only solutions
\emph{spanning} $S$ and its complement --- or spanning two distinct sources inside $S$
--- are lost. This is the narrower fact the impossibility argument needs.
\end{remark}

\begin{remark}[Side conditions]
The construction needs $B'$ connected (two patterns sharing a variable), $S \neq
\emptyset$, and at least one pattern of $B'$ matched within $S$; otherwise there is
nothing for the adversary's triple to join to. Taking $B' = B$ removes any further
condition, since then no join partner in $B \setminus B'$ is required.
\end{remark}

\begin{remark}[What this asks of the assumption]
The argument turns entirely on the adversary being free to place arbitrary data behind an
unfetched URI. Any fix must remove that freedom rather than improve the source
descriptions: no refinement of $\Delta$ over $\Dd'(T)$ helps, because both worlds agree
there. This is what coverage and authority declarations supply in the next section.
\end{remark}

\section{Coverage}

\begin{definition}[Coverage]
$\mathrm{cov} : \Dd \to 2^{\Dd}$, in two flavours:
\[
  d \in \mathrm{cov}_{\mathsf{agg}}(c) \iff \dataf(d) \subseteq \dataf(c),
  \qquad
  d \in \mathrm{cov}_{\mathsf{idx}}(x) \iff x \text{ publishes } \Delta(\cdot,d).
\]
An aggregator lets you \emph{act}; an index lets you \emph{decide}. Lift to sets by
$\mathrm{cov}(A) := \bigcup_{c \in A}\mathrm{cov}(c)$, and call $A$ a \emph{cover} of
$\Dd$ if $\mathrm{cov}(A) = \Dd$.
\end{definition}

\begin{remark}
$\mathrm{cov}$ is extensive, monotone and idempotent (coverage is transitive), hence a
closure operator on $2^{\Dd}$. Covers form an up-set; the minimal ones are the objects
of interest.
\end{remark}

\begin{lemma}[Answer equivalence]
If $A$ is an aggregator-cover of $\Dd$ then $G_A = G_{\Dd}$, and therefore
$\ev{Q}{G_A} = \ev{Q}{G_{\Dd}}$ for every monotone $Q$.
\end{lemma}

\begin{lemma}[Observation asymmetry]
Let $\Dd'_t(T) := \{\, d \in \Dd'(T) \mid \Delta(t,d) = 1 \,\} \subseteq \Dd_t$. Then
\[
  |\Dd'_t(T)| \geq 2 \;\Longrightarrow\; |\Dd_t| \geq 2,
  \qquad\text{but}\qquad
  |\Dd'_t(T)| = 1 \;\not\Longrightarrow\; |\Dd_t| = 1 .
\]
\end{lemma}

\noindent
Non-exclusivity is soundly concludable from a finite observation; exclusivity is not.
That asymmetry \emph{is} the impossibility result.

\begin{definition}[Discovery closure]
Recall $\disc(d, \dataf(d))$ from Definition~\ref{def:setting}. A set $A$ is
\emph{discovery-closed} if
\[
  \forall c \in A : \ \disc\bigl(c, \dataf(c)\bigr) \subseteq \mathrm{cov}(A).
\]
\end{definition}

\begin{lemma}[Frontier]
If $A$ is discovery-closed and $\mathsf{seeds} \subseteq A$, then
$\mathrm{reach}(\mathsf{seeds}) \subseteq \mathrm{cov}(A)$.
\end{lemma}

\noindent
This converts the unverifiable global condition ``I have seen everything'' into the
local, checkable one ``my frontier is covered''.

\begin{theorem}[Decidability under coverage]
If at time $T$ some $A \subseteq \Dd'(T)$ is discovery-closed with admissible index
coverage, then $\Dd_t$ is computable from $\Dd'(T)$ alone, and exclusivity of $k$ for
$\{t_1,\dots,t_n\}$ is decidable at time $T$.
\end{theorem}

\begin{theorem}[Impossibility without coverage]
%% specialisation of Theorem~\ref{thm:nocert} to the coverage setting
If $\mathrm{cov}(\Dd'(T)) \subsetneq \Dd$, then for every $k$ with $k \in \Dd_{t_i}$ for
all $i$ there is a completion of $\Dd$ consistent with every observation made up to $T$
under which delegating $\{t_1,\dots,t_n\}$ to $k$ is incorrect.
\end{theorem}

\begin{corollary}
Without coverage declarations, exclusivity is \emph{eventually computable} but not
\emph{finitely computable} (Hartig \& Freytag, Defs.~13--14). With an admissible
discovery-closed index cover it becomes finitely computable.
\end{corollary}

\begin{remark}
Choosing the cheapest $A \subseteq \Dd'(T)$ with $\mathrm{cov}(A) = \Dd$ is weighted Set
Cover: NP-hard, greedy gives $H_n \approx \ln n$. FedQPL already proves the analogous
hardness for minimal source assignments by reduction from node cover (Def.~10 and
Prop.~2), so that precedent is citable rather than new.
\end{remark}

\begin{remark}
Coverage and the namespace scope claim are one mechanism: a claim over $N$ covers
exactly the sources whose data lies in $S_N$. Coverage is the extensional presentation,
the scope claim the intensional one.
\end{remark}

\end{document}
